%!TEX root = ../thesis.tex
\chapter{Introduction}
\label{ch:intro}
\todo{sämtliche Abkürzungen mit glocaries einbinden und verwalten}
\section{Motivation}
%Proposal Version
    % Durch die leichte Verfügbarkeit von hochauflösenden Satelliten- und Luftbildern ist Detektion von Objekten weltweit möglich. Da außerdem die hohe Rechenleistung, die für maschinelles Lernen erforderlich ist, immer preiswerter wird, wird die Anwendung von künstlicher Intelligenz zur vollautomatischen Analyse immer leichter. \\
    % Das Monitoring von zerstörten (militärischen) Gerät oder Gebäudestrukturen in Konfliktregionen, wie Darfur \cite{Knoth2017} oder der Ukraine, kann durch künstliche Intelligenz unterstützt werden. Hier können Luftbiler gut genutzt werden, weil das Erstellen von Fotos vor Ort durch die Sicherheitslage am Boden in Konfliktregionen gefährlich sein kann. \\
    % Das Monitoring von militärischem Gerät an Ländergrenzen könnte benutzt werden um eine genauere Lageeinschätzung zu generieren, ob ein Konflikt bevorsteht. Es ist auch möglich die Arbeit von menschlichen Analysten unterstützen, da die Modelle Vorschläge für mögliche Detektionen und Klassifizierungen von Objekten liefern. \\
    % Maschinelles Lernen kann zur Automatisierung angewendet werden, um größere Datensätze zu verarbeiten. Eine Herausforderung hier besteht im Training der Algorithmen, da die Verfügbarkeit von Trainingsdaten limitiert ist.\\

    % Ein Anwendungszenario wäre die Auswertung des menschlichen Mobilitätsverhaltens im Rahmen des Zählen von Autos auf großen öffentlichen Parkplätzen. Hier kann auch ein Tracking von Verkehrsaufkommen durch die KI landesweit erfolgen, wenn Fahrzeuge auf einem Satellitenbild detektiert werden können. Dieser Ansatz kann genutzt werden um abzuschätzen ob mehr Parkflächen erforderlich sind oder die Anzahl der aktuellen Flächen ausreicht. Weitere Fragestellungen wären, ob es möglich ist fahrende und stehende Autos zu unterscheiden.

    % Außerdem kann in dieser Arbeit evaluiert werden, ob die Präzision der  Objekterkennung und -klassifzierung von Fahrzeugen in elektro-optischen, multi-spektralen Luftaufnahmen durch die Nutzung von mehr Bildkanälen (NIR, IR) verbessert werden kann. Dies kann durch einen Vergleich der Ergebnisse beim Trainieren desselben Deep Learning Modells geschehen, welches jeweils nur mit 3 und einmal mit  mehr Kanälen trainiert wurde. Da Fahrzeuge relativ klein auf hochauflösenden Satellitenbildern dargestellt werden, kann auch evaluiert werden, wie gut sich ein Deep Learning Modell zur Klassifizierung sehr kleiner Objekte eignet, bzw. ob ein weiterer Kanal die Präzision verbessert.
    % \todo{kann yolov9 kleine objeket auf satbildern zuverlässig erkennen? Vergleich mit Josis Thesis?}

%Ohne Josis Thesis


% Durch die zunehmende Verfügbarkeit hochauflösender Satelliten- und Luftbilddaten ist die weltweite Detektion von Objekten technisch möglich geworden. Gleichzeitig sinken die Kosten für die hohe Rechenleistung, die für maschinelles Lernen erforderlich ist, sodass der Einsatz von Methoden der Künstlichen Intelligenz (KI) zur vollautomatischen Analyse solcher Daten kontinuierlich erleichtert wird. Damit entstehen neue Möglichkeiten zur effizienten Auswertung großer Mengen an Bilddaten, die sowohl im zivilen als auch im militärischen Kontext von erheblicher Relevanz sind. \\

% Ein bedeutsamer Anwendungsbereich liegt im Monitoring von zerstörten militärischen Anlagen oder Gebäudestrukturen in Konfliktregionen, wie beispielsweise in Darfur \cite{Knoth2017} oder der Ukraine. Der Einsatz von KI-gestützter Bildanalyse kann hier wertvolle Unterstützung bieten, da eine direkte fotografische Dokumentation vor Ort aufgrund der unsicheren Sicherheitslage in Konfliktgebieten oftmals nicht möglich oder mit erheblichen Risiken verbunden ist. Luft- und Satellitenbilder eignen sich daher in besonderem Maße für die Informationsgewinnung. \\

% Darüber hinaus könnte die kontinuierliche Überwachung von militärischem Gerät an Ländergrenzen zur Verbesserung der Lageeinschätzung beitragen, indem potenzielle Eskalationen oder bevorstehende Konflikte frühzeitig erkannt werden. Ein entscheidender Vorteil der KI-gestützten Analyse besteht hierbei in der Möglichkeit, menschliche Analysten bei ihrer Arbeit zu entlasten, indem Modelle Vorschläge für Detektionen und Klassifizierungen liefern, die anschließend überprüft und validiert werden können. Auf diese Weise können bestehende Arbeitsabläufe effizienter gestaltet und Entscheidungsträger schneller mit relevanten Informationen versorgt werden. \\

% Neben diesen sicherheitsrelevanten Fragestellungen existieren auch vielfältige zivile Anwendungsbereiche. So kann die Analyse des Mobilitätsverhaltens der Bevölkerung durch die Erfassung von Fahrzeugen auf großen öffentlichen Parkplätzen unterstützt werden. Mithilfe automatisierter Detektions- und Klassifikationsmethoden lassen sich Fahrzeuge auf Satelliten- oder Luftbildern identifizieren und zählen, was wiederum Rückschlüsse auf die Auslastung von Parkflächen ermöglicht. Diese Informationen können beispielsweise genutzt werden, um den Bedarf an zusätzlichen Parkflächen zu bestimmen oder die Effizienz der bestehenden Infrastruktur zu evaluieren. Darüber hinaus eröffnet sich die Möglichkeit, das Verkehrsaufkommen auf nationaler Ebene zu quantifizieren, sofern eine kontinuierliche Fahrzeugdetektion auf Satellitenbildern implementiert wird. Eine weiterführende Fragestellung ist die Differenzierung zwischen fahrenden und stehenden Fahrzeugen, die zusätzliche Einblicke in Mobilitätsmuster und Verkehrsströme liefern könnte. \\

% Eine zentrale Herausforderung bei der Anwendung maschinellen Lernens auf hochauflösende Fernerkundungsbilder liegt jedoch in der begrenzten Verfügbarkeit geeigneter Trainingsdaten sowie in der Tatsache, dass Fahrzeuge auf Satellitenaufnahmen aufgrund des Maßstabs relativ klein dargestellt werden. Die präzise Erkennung und Klassifizierung dieser kleinen Objekte erfordert angepasste Deep-Learning-Methoden, die in der Lage sind, auch feine Strukturen zuverlässig zu erfassen. \\

% Im Rahmen dieser Arbeit soll daher untersucht werden, inwieweit die Präzision der Objekterkennung und -klassifizierung von Fahrzeugen in elektro-optischen und multispektralen Luft- und Satellitenaufnahmen verbessert werden kann, wenn zusätzliche Bildkanäle wie beispielsweise NIR- oder IR-Bänder einbezogen werden. Hierfür wird ein Deep-Learning-Modell sowohl mit konventionellen dreikanaligen (RGB) Daten als auch mit multispektralen Daten trainiert und die Ergebnisse miteinander verglichen. Auf diese Weise kann evaluiert werden, ob die Erweiterung der Spektralinformation die Erkennungsleistung für kleine Objekte wie Fahrzeuge verbessert. Gleichzeitig wird überprüft, wie gut sich moderne Deep-Learning-Ansätze grundsätzlich zur Detektion und Klassifizierung sehr kleiner Objekte eignen und welche Potenziale sich durch den Einsatz multispektraler Fernerkundungsdaten eröffnen.

%ChatGPT Vorschlag mit Josis Thesis
% Durch die zunehmende Verfügbarkeit hochauflösender Satelliten- und Luftbilddaten ist die weltweite Detektion von Objekten technisch möglich geworden. Gleichzeitig sinken die Kosten für die hohe Rechenleistung, die für maschinelles Lernen erforderlich ist, sodass der Einsatz von Methoden der Künstlichen Intelligenz (KI) zur vollautomatischen Analyse solcher Daten kontinuierlich erleichtert wird. Damit entstehen neue Möglichkeiten zur effizienten Auswertung großer Mengen an Bilddaten, die sowohl im zivilen als auch im militärischen Kontext von erheblicher Relevanz sind. \\

% Ein bedeutsamer Anwendungsbereich liegt im Monitoring von zerstörten militärischen Anlagen oder Gebäudestrukturen in Konfliktregionen, wie beispielsweise in Darfur \cite{Knoth2017} oder der Ukraine. Der Einsatz von KI-gestützter Bildanalyse kann hier wertvolle Unterstützung bieten, da eine direkte fotografische Dokumentation vor Ort aufgrund der unsicheren Sicherheitslage in Konfliktgebieten oftmals nicht möglich oder mit erheblichen Risiken verbunden ist. Luft- und Satellitenbilder eignen sich daher in besonderem Maße für die Informationsgewinnung. \\

% Darüber hinaus könnte die kontinuierliche Überwachung von militärischem Gerät an Ländergrenzen zur Verbesserung der Lageeinschätzung beitragen, indem potenzielle Eskalationen oder bevorstehende Konflikte frühzeitig erkannt werden. Ein entscheidender Vorteil der KI-gestützten Analyse besteht hierbei in der Möglichkeit, menschliche Analysten bei ihrer Arbeit zu entlasten, indem Modelle Vorschläge für Detektionen und Klassifizierungen liefern, die anschließend überprüft und validiert werden können. Auf diese Weise können bestehende Arbeitsabläufe effizienter gestaltet und Entscheidungsträger schneller mit relevanten Informationen versorgt werden. \\

% Neben diesen sicherheitsrelevanten Fragestellungen existieren auch vielfältige zivile Anwendungsbereiche. So kann die Analyse des Mobilitätsverhaltens der Bevölkerung durch die Erfassung von Fahrzeugen auf großen öffentlichen Parkplätzen unterstützt werden. Mithilfe automatisierter Detektions- und Klassifikationsmethoden lassen sich Fahrzeuge auf Satelliten- oder Luftbildern identifizieren und zählen, was wiederum Rückschlüsse auf die Auslastung von Parkflächen ermöglicht. Diese Informationen können beispielsweise genutzt werden, um den Bedarf an zusätzlichen Parkflächen zu bestimmen oder die Effizienz der bestehenden Infrastruktur zu evaluieren. Darüber hinaus eröffnet sich die Möglichkeit, das Verkehrsaufkommen auf nationaler Ebene zu quantifizieren, sofern eine kontinuierliche Fahrzeugdetektion auf Satellitenbildern implementiert wird. Eine weiterführende Fragestellung ist die Differenzierung zwischen fahrenden und stehenden Fahrzeugen, die zusätzliche Einblicke in Mobilitätsmuster und Verkehrsströme liefern könnte. \\

% Eine zentrale Herausforderung bei der Anwendung maschinellen Lernens auf hochauflösende Fernerkundungsbilder liegt jedoch in der begrenzten Verfügbarkeit geeigneter Trainingsdaten sowie in der Tatsache, dass Fahrzeuge auf Satellitenaufnahmen aufgrund des Maßstabs relativ klein dargestellt werden. Die präzise Erkennung und Klassifizierung dieser kleinen Objekte erfordert angepasste Deep-Learning-Methoden, die in der Lage sind, auch feine Strukturen zuverlässig zu erfassen. \\

% Vorarbeiten in diesem Themenfeld haben bereits gezeigt, dass Deep-Learning-Methoden wie YOLO grundsätzlich in der Lage sind, kleine Objekte in hochauflösenden Fernerkundungsbildern zu detektieren. Dabei konnte nachgewiesen werden, dass durch gezielte Anpassung von Hyperparametern wie Rastergröße, Trainingsdauer und Learning Rate die mittlere Average Precision (mAP) für sehr kleine Objekte verbessert werden kann. Insbesondere beim Airbus Ship Detection Datensatz wurden durch diese Optimierungen deutliche Leistungssteigerungen in den Klassen kleiner Objekte erzielt. \\

% Allerdings zeigten dieselben Untersuchungen auch die Grenzen dieses Ansatzes auf: Für komplexere Datensätze mit größeren Bildausschnitten und einer Vielzahl verschiedener Objektklassen, wie etwa den DOTA-Datensatz, ließ sich keine signifikante Verbesserung der Erkennungsleistung erzielen. Daraus ergibt sich die Schlussfolgerung, dass die Nutzung von YOLO für komplexe Szenarien mit mehreren Objektklassen und sehr kleinen Zielobjekten nicht zielführend ist. \\

%Mit Josis THesis

% Durch die zunehmende Verfügbarkeit hochauflösender Satelliten- und Luftbilddaten ist die weltweite Detektion von Objekten technisch möglich geworden. Gleichzeitig sinken die Kosten für die erforderliche Rechenleistung, sodass der Einsatz von Methoden der Künstlichen Intelligenz (KI) zur vollautomatischen Analyse solcher Daten kontinuierlich erleichtert wird. Damit eröffnen sich neue Möglichkeiten zur effizienten Auswertung großer Mengen an Bilddaten, die sowohl im zivilen als auch im militärischen Kontext von erheblicher Relevanz sind. \\

% Ein bedeutender Anwendungsbereich liegt in der Analyse und Dokumentation von zerstörten Infrastrukturen oder Gebäudestrukturen in Krisen- und Katastrophengebieten, beispielsweise wenn eine direkte Erhebung vor Ort aufgrund von Sicherheitsrisiken nicht möglich ist \cite{Knoth2017}. KI-gestützte Bildanalyse kann hier wertvolle Unterstützung leisten, da Luft- und Satellitenbilder eine effiziente und sichere Informationsgewinnung ermöglichen.
% \begin{itemize}
%     \item Bewaffnete Konflikte haben einen Einfluss auf die Ernährungssicherheit und nachhaltige Bewirtschaftungen von Land- und Wasseressourcen \cite{Mhanna2023}
%     \item Herausforderung ist der Mangel an direkt gemessenden Ground Truth Daten bspw im sysrischen Teil des grenzüberschreitenden Orontes Flussbecken seit 2013; aufgrund der sicherheitlichen Situation
%     \item es wurden Methodiken entwickelt, um Land Used Land Cover Veränderungen in dem Gebiet zwischen 2004 und 2022 besser zu verstehen
%     \item Anbauflächen gingen stark zurück und die Entwicklung von FLüchtlingssiedlugnen konnte verfolgt werden \cite{Mhanna2023}
% \end{itemize}

% \begin{itemize}
%     \item Weitere Ansätze umfassen die Kartierung von Umweltauswirkungen und Landnarben durch die weltweit zunehmende Anzahl von Kriegen und bewaffneten Konfliktgebieten
%     \item OBIA wurde für diesen Zweck mithilfe von DL CNNs am Beispiel des Iran-Irak Krieges angewendet
%     \item Übertragbarkeit auf den syrischen Bürgerkrieg und den Karabach Konflikt (Aserbaidschan-Armenien) untersucht
%     \item Genauigkeit von 0,96 im Iran Irak Krieg und 0.93 bzw. 09.4 für Landnarben in Syrien und Karabach zur Erkennung von Landnarben \cite{Feizizadeh2025}
% \end{itemize}

% \begin{itemize}
%     \item detektion von massengräbern ist aufgrund der sicherheitssituation am boden ebenfalls oft sehr erschwert
%     \item in der vergangneheit stützte man sich auf Zeugenaussagen
%     \item andere Ortungsmehtoden setzen meist voraus, dass der Standrot des Grabes ziemlich genau bekannt ist und das Untersuchungsgebiet physische betreten werde nkann
%     \item Fernerkundung kann deutlich größere Gebiete abdecken ohne das Untersuchungspersonal Risiken auszusetzen, oder durch politische Akteure an den Untersuchungne gehindert werde nkann
%     \item Langzeit-Experimente zeigen, dass Hyperspektralbilder mithilfe von Fernerkundung ein Leistungsfähiges Detektionsinstrument ist um Massengräber in Echtzeit oder in bestimmten Fällen sogar Rückwirketnd von seinr umgebung zu unterscheidne \cite{Kalacska2006}
%     \item gleiches wurde auch für einezelne gräber erprbot und untersucht \cite{Leblanc2014}, obwohl dies aufgrund der geringeren Größe und Masse der Ziele deutlich schwieriger ist. ergebneisse zeigetn, dass hyperspektrale Bildgebung vielversprechend für die Erkennung ovn begrabenen Überresten ist \cite{Leblanc2014}
% \end{itemize}

% \begin{itemize}
%     \item Die genannten Anwendungsfälle können zur Beweisführugn und Nachverfolgung von (Kriegs-)Verbrechen genutzt werden um die Verwortlichen zur Rechenschaft zu ziehen und auch um Vermisste zu identifizieren
%     \item Um Konflikte im Vorhinein zu verhindern, bzw. die Anbahnugn militärischer Konflikte zu erkennen kann das Mobilitätsmuster über große Gebiete hinweg analysiert werden
% \end{itemize}

% Es kann die Analyse des Mobilitätsverhaltens der Bevölkerung durch die Erfassung von Fahrzeugen auf großen öffentlichen Parkplätzen unterstützt werden. Mithilfe automatisierter Detektions- und Klassifikationsmethoden lassen sich Fahrzeuge auf Satelliten- oder Luftbildern identifizieren und zählen, was Rückschlüsse auf die Auslastung von Parkflächen erlaubt. Diese Informationen können beispielsweise genutzt werden, um den Bedarf an zusätzlichen Parkflächen zu bestimmen oder die Effizienz der bestehenden Infrastruktur zu evaluieren. Zudem eröffnet sich die Möglichkeit, das Verkehrsaufkommen auf regionaler oder nationaler Ebene zu quantifizieren, sofern eine kontinuierliche Fahrzeugdetektion implementiert wird. Eine weiterführende Fragestellung ist die Differenzierung zwischen fahrenden und stehenden Fahrzeugen, die zusätzliche Einblicke in Mobilitätsmuster und Verkehrsströme liefern kann.

% Auch innerhalb sicherheitsrelevanter Anwendungen eröffnen sich vielfältige Einsatzmöglichkeiten. Die automatisierte Auswertung großer Bilddatenmengen kann dazu beitragen, Veränderungen in sensiblen Regionen frühzeitig zu erkennen und die Lageeinschätzung zu verbessern. Ein wesentlicher Vorteil besteht darin, dass durch KI-Modelle menschliche Analysten entlastet werden, indem potenzielle Objekte oder Strukturen automatisch erkannt und markiert werden, bevor eine abschließende Validierung erfolgt.\\

% Eine zentrale Herausforderung bei der Anwendung maschinellen Lernens auf hochauflösende Fernerkundungsbilder liegt in der begrenzten Verfügbarkeit geeigneter Trainingsdaten sowie darin, dass Fahrzeuge aufgrund des Maßstabs relativ klein dargestellt werden. Die präzise Erkennung und Klassifizierung dieser kleinen Objekte erfordert angepasste Deep-Learning-Methoden, die auch feine Strukturen zuverlässig erfassen können. \\

% Vorarbeiten in diesem Themenfeld, wie sie in einer vorherigen Bachelorarbeit \cite{Balzer2022} durchgeführt wurden, zeigen, dass Deep-Learning-Methoden wie \acrfull{YOLO} (v3) kleine Objekte in einfachen Szenarien grundsätzlich erkennen können. Dabei konnte durch die Anpassung von Hyperparametern die \acrfull{mAP} für sehr kleine Objekte verbessert werden. Allerdings wird im Anwendungsfall der Schiffsdetektion deutlich, dass die erzielte mAP kritisch zu bewerten ist und die Eignung von YOLOv3 für eine zuverlässige Detektion von kleinen Objekten nicht uneingeschränkt gegeben ist \cite{Balzer2022}. \\

% Die vorliegende Masterarbeit knüpft an diese Ergebnisse an und überträgt die Fragestellung auf den Bereich der Detektion und Klassifikation von kleinen Fahrzeugen in hochauflösenden multispektralen Luftbilddaten. Dabei wird eine neuere Version von \acrfull{YOLO} (v9) eingesetzt, die Verbesserungen in der Netzwerkarchitektur aufweist. Zusätzlich wird untersucht, inwieweit die Integration elektro-optische und multispektrale Luftbildaufnahmen  die Präzision bei der Detektion und Klassifikation sehr kleiner Objekte steigern kann. Auf diese Weise erweitert die Arbeit bestehende Forschung sowohl inhaltlich (Fokus auf Fahrzeuge) als auch methodisch (Verwendung moderner YOLO-Architektur und multispektraler Daten) und adressiert die Forschungslücke der Objekterkennung in komplexen Szenarien.

%ohne militärischen kontext

% Durch die zunehmende Verfügbarkeit hochauflösender Satelliten- und Luftbilddaten ist die weltweite Detektion von Objekten technisch möglich geworden. Gleichzeitig sinken die Kosten für die erforderliche Rechenleistung, sodass der Einsatz von Methoden der Künstlichen Intelligenz (KI) zur vollautomatischen Analyse solcher Daten kontinuierlich erleichtert wird. Damit eröffnen sich neue Möglichkeiten zur effizienten Auswertung großer Mengen an Bilddaten, die in verschiedenen zivilen Bereichen von erheblicher Relevanz sind. \\

% Ein bedeutsamer Anwendungsbereich liegt im Monitoring von Gebäudestrukturen oder städtischen Infrastrukturen. Der Einsatz von KI-gestützter Bildanalyse kann hier wertvolle Unterstützung bieten, insbesondere wenn eine direkte Datenerhebung vor Ort schwierig oder mit hohem Aufwand verbunden ist. Luft- und Satellitenbilder eignen sich daher besonders gut zur Informationsgewinnung. \\

% Auch im zivilen Bereich eröffnen sich vielfältige Anwendungsfelder. So kann die Analyse des Mobilitätsverhaltens der Bevölkerung durch die Erfassung von Fahrzeugen auf großen öffentlichen Parkplätzen unterstützt werden. Mithilfe automatisierter Detektions- und Klassifikationsmethoden lassen sich Fahrzeuge auf Satelliten- oder Luftbildern identifizieren und zählen, was Rückschlüsse auf die Auslastung von Parkflächen erlaubt. Diese Informationen können beispielsweise genutzt werden, um den Bedarf an zusätzlichen Parkflächen zu bestimmen oder die Effizienz der bestehenden Infrastruktur zu evaluieren. Zudem eröffnet sich die Möglichkeit, das Verkehrsaufkommen auf regionaler oder nationaler Ebene zu quantifizieren, sofern eine kontinuierliche Fahrzeugdetektion implementiert wird. Eine weiterführende Fragestellung ist die Differenzierung zwischen fahrenden und stehenden Fahrzeugen, die zusätzliche Einblicke in Mobilitätsmuster und Verkehrsströme liefern kann. \\

% Eine zentrale Herausforderung bei der Anwendung maschinellen Lernens auf hochauflösende Fernerkundungsbilder liegt in der begrenzten Verfügbarkeit geeigneter Trainingsdaten sowie darin, dass Fahrzeuge aufgrund des Maßstabs relativ klein dargestellt werden. Die präzise Erkennung und Klassifizierung dieser kleinen Objekte erfordert angepasste Deep-Learning-Methoden, die auch feine Strukturen zuverlässig erfassen können. \\

% Vorarbeiten in diesem Themenfeld, wie sie in einer vorherigen Bachelorarbeit durchgeführt wurden, zeigen, dass Deep-Learning-Methoden wie \acrfull{YOLO} (v3) kleine Objekte in einfachen Szenarien grundsätzlich erkennen können. Dabei konnte durch die Anpassung von Hyperparametern die \acrfull{mAP} für sehr kleine Objekte verbessert werden. Allerdings wird im Anwendungsfall der Fahrzeugdetektion deutlich, dass die erzielte mAP kritisch zu bewerten ist und die Eignung von YOLOv3 für eine zuverlässige Detektion von kleinen Objekten nicht uneingeschränkt gegeben ist \cite{Balzer2022}. \\

% Die vorliegende Masterarbeit knüpft an diese Ergebnisse an und überträgt die Fragestellung auf den Bereich der Detektion und Klassifikation von kleinen Fahrzeugen in hochauflösenden multispektralen Luftbilddaten. Dabei wird eine neuere Version von \acrfull{YOLO} (v9) eingesetzt, die Verbesserungen in der Netzwerkarchitektur aufweist. Zusätzlich wird untersucht, inwieweit die Integration elektro-optischer und multispektraler Luftbildaufnahmen die Präzision bei der Detektion und Klassifikation sehr kleiner Objekte steigern kann. Auf diese Weise erweitert die Arbeit bestehende Forschung sowohl inhaltlich (Fokus auf Fahrzeuge) als auch methodisch (Verwendung moderner YOLO-Architektur und multispektraler Daten) und adressiert die Forschungslücke der Objekterkennung in komplexen Szenarien.

Durch die zunehmende Verfügbarkeit hochauflösender Satelliten- und Luftbilddaten ist die weltweite Detektion von Objekten technisch möglich geworden. Gleichzeitig sinken die Kosten für die erforderliche Rechenleistung, sodass der Einsatz von Methoden der Künstlichen Intelligenz (KI) zur vollautomatischen Analyse solcher Daten kontinuierlich erleichtert wird. Damit eröffnen sich neue Möglichkeiten zur effizienten Auswertung großer Mengen an Bilddaten, die sowohl im zivilen als auch im militärischen Kontext von erheblicher Relevanz sind \cite{Lemley2017}. 

Ein bedeutender Anwendungsbereich liegt in der Analyse und Dokumentation zerstörter Infrastrukturen oder Gebäudestrukturen in Krisen- und Katastrophengebieten, insbesondere wenn eine direkte Erhebung vor Ort aufgrund von Sicherheitsrisiken nicht möglich ist \cite{Knoth2017}. KI-gestützte Bildanalyse kann hier wertvolle Unterstützung leisten, da Luft- und Satellitenbilder eine effiziente und sichere Informationsgewinnung ermöglichen.

Bewaffnete Konflikte haben tiefgreifende Auswirkungen auf die Ernährungssicherheit sowie auf die nachhaltige Bewirtschaftung von Land- und Wasserressourcen. Durch Zerstörung von Infrastruktur, Fluchtbewegungen und den Verlust agrarischer Produktionsflächen entstehen komplexe sozioökologische Krisen, die die Lebensgrundlagen der lokalen Bevölkerung unmittelbar gefährden. Studien zeigen, dass insbesondere in konfliktbetroffenen Regionen wie dem Nahen Osten die Resilienz landwirtschaftlicher Systeme erheblich beeinträchtigt wird \cite{Mhanna2023}. Die langfristigen Folgen dieser Entwicklungen wirken sich nicht nur auf die regionale, sondern auch auf die globale Ernährungssicherheit aus.

Eine zentrale Herausforderung in der wissenschaftlichen Analyse solcher Gebiete besteht im Mangel an direkt gemessenen \acrfull{GT}-Daten. Im syrischen Teil des grenzüberschreitenden Orontes-Flussbeckens beispielsweise sind seit 2013 aufgrund der sicherheitspolitischen Lage keine systematischen Feldmessungen mehr möglich \cite{Mhanna2023}. Der eingeschränkte Zugang zu den Untersuchungsgebieten erschwert die Verifizierung satellitengestützter Fernerkundungsdaten erheblich und stellt somit ein wesentliches methodisches Hindernis für die Forschung dar.

Um dennoch Veränderungen der Landnutzung und Landbedeckung (\acrfull{LULC}) in Konfliktregionen zu erfassen, wurden verschiedene methodische Ansätze entwickelt. Für den Zeitraum von 2004 bis 2022 konnten mit Hilfe multitemporaler Satellitenbildanalysen signifikante Muster identifiziert werden, die ein besseres Verständnis der ökologischen und sozioökonomischen Dynamiken im Orontes-Becken ermöglichen. Dabei wurde insbesondere deutlich, dass landwirtschaftlich genutzte Flächen drastisch zurückgingen, während gleichzeitig die Entstehung und Ausdehnung von Flüchtlingssiedlungen nachvollzogen werden konnte \cite{Mhanna2023}.

Darüber hinaus wurden in der jüngeren Forschung zunehmend Ansätze entwickelt, die sich mit der Kartierung von Umweltauswirkungen und sogenannten \emph{Landnarben} infolge bewaffneter Konflikte beschäftigen. Die wachsende Zahl kriegerischer Auseinandersetzungen weltweit macht es erforderlich, die durch militärische Aktivitäten verursachten landschaftlichen Veränderungen systematisch zu erfassen. Hierfür wurde die objektbasierte Bildanalyse (\acrfull{OBIA}) unter Einsatz tiefer neuronaler Netze (\acrfull{DCNN}) angewandt, beispielsweise im Kontext des Iran-Irak-Krieges. Die erzielten Ergebnisse zeigen eine hohe Genauigkeit bei der Erkennung konfliktbedingter Landschaftsveränderungen. Zudem konnte die Übertragbarkeit der entwickelten Methoden auf andere Konfliktszenarien, wie den syrischen Bürgerkrieg und den Karabach-Konflikt zwischen Aserbaidschan und Armenien, erfolgreich untersucht werden \cite{Feizizadeh2025}.

Ein weiteres sensibles, jedoch hochrelevantes Forschungsfeld betrifft die Detektion von Massengräbern. Aufgrund der gefährlichen Sicherheitslage ist eine direkte Untersuchung am Boden in vielen Konfliktgebieten kaum möglich \cite{Kalacska2006}. Historisch stützte sich die Lokalisierung solcher Stätten auf Zeugenaussagen oder auf Methoden, die einen bekannten Standort und physischen Zugang voraussetzen. Fernerkundung hingegen bietet die Möglichkeit, großflächige Gebiete systematisch und sicher zu analysieren, ohne dass Untersuchungspersonal einem Risiko ausgesetzt wird oder politische Akteure den Zugang behindern können. Langzeitstudien zeigen, dass hyperspektrale Bildgebung ein leistungsfähiges Instrument zur Detektion von Massengräbern darstellt, da sie spektrale Anomalien in der Vegetation und im Boden erkennt, die auf Zersetzungsprozesse hinweisen \cite{Kalacska2006}. Diese Technologie wurde auch erfolgreich zur Identifikation einzelner Gräber getestet, wenngleich dies aufgrund der geringeren Größe und Masse der Objekte methodisch anspruchsvoller ist. Die Ergebnisse deuten jedoch darauf hin, dass hyperspektrale Fernerkundung ein vielversprechendes Werkzeug zur Erkennung begrabener Überreste darstellt \cite{Leblanc2014}.

Ein weiteres Anwendungsszenario stellt die Analyse des Mobilitätsverhaltens der Bevölkerung dar, beispielsweise durch die Erfassung von Fahrzeugen auf großen öffentlichen Parkplätzen. Mithilfe automatisierter Detektions- und Klassifikationsmethoden lassen sich Fahrzeuge auf Satelliten- oder Luftbildern identifizieren und zählen, was Rückschlüsse auf die Auslastung von Parkflächen erlaubt. Diese Informationen können genutzt werden, um den Bedarf an zusätzlichen Parkflächen zu bestimmen oder die Effizienz bestehender Infrastrukturen zu evaluieren. Zudem eröffnet sich die Möglichkeit, das Verkehrsaufkommen auf regionaler oder nationaler Ebene zu quantifizieren, sofern eine kontinuierliche Fahrzeugdetektion implementiert wird \cite{Eikvil2009}. Eine weiterführende Fragestellung ist die Differenzierung zwischen fahrenden und stehenden Fahrzeugen, die zusätzliche Einblicke in Mobilitätsmuster und Verkehrsströme liefern kann.

Auch im sicherheitsrelevanten Kontext ergeben sich vielfältige Einsatzmöglichkeiten. Die automatisierte Auswertung großer Bilddatenmengen kann dazu beitragen, Veränderungen in sensiblen Regionen frühzeitig zu erkennen und die Lageeinschätzung zu verbessern. Ein wesentlicher Vorteil besteht darin, dass durch KI-Modelle menschliche Analysten entlastet werden, indem potenzielle Objekte oder Strukturen automatisch erkannt und markiert werden, bevor eine abschließende Validierung erfolgt.

Die in den vorangegangenen Abschnitten dargestellten Anwendungsfälle verdeutlichen das Potenzial der Fernerkundung für humanitäre, ökologische und forensische Fragestellungen. Die erfassten Daten können zur Beweisführung und Nachverfolgung von Kriegsverbrechen herangezogen werden, um Verantwortliche zur Rechenschaft zu ziehen und Vermisste zu identifizieren. Darüber hinaus eröffnen die gewonnenen Erkenntnisse Möglichkeiten zur Früherkennung und Prävention von Konflikten. Durch die Analyse von Mobilitätsmustern und Veränderungen in der Landnutzung über große Gebiete hinweg lassen sich potenzielle Spannungsfelder frühzeitig erkennen und geeignete Maßnahmen zur Konfliktminderung ableiten.

Eine zentrale Herausforderung bei der Anwendung maschinellen Lernens auf hochauflösende Fernerkundungsbilder liegt jedoch in der begrenzten Verfügbarkeit geeigneter Trainingsdaten sowie darin, dass Fahrzeuge aufgrund des Maßstabs relativ klein dargestellt werden. Die präzise Erkennung und Klassifizierung dieser kleinen Objekte erfordert spezialisierte Deep-Learning-Modelle, die auch feine Strukturen zuverlässig erfassen können.

Vorarbeiten in diesem Themenfeld, wie sie in einer vorherigen Bachelorarbeit \cite{Balzer2022} durchgeführt wurden, zeigen, dass Deep-Learning-Methoden wie \acrfull{YOLO} (v3) kleine Objekte in einfachen Szenarien grundsätzlich erkennen können. Durch gezielte Anpassung von Hyperparametern konnte die \acrfull{mAP} für sehr kleine Objekte verbessert werden. Allerdings zeigte sich im Anwendungsfall der Schiffsdetektion, dass die erzielten mAP-Werte kritisch zu bewerten sind und die Eignung von YOLOv3 für die zuverlässige Detektion kleiner Objekte nur eingeschränkt gegeben ist \cite{Balzer2022}.

Deep Learning eignet sich besonders für die Analyse von Fernerkundungsdaten, da es komplexe räumlich-semantische Muster in hochauflösenden Bildern erfassen kann, die mit klassischen, feature-basierten Methoden nur schwer zu modellieren sind. Insbesondere die Kombination aus multispektralen Kanälen, wie \acrshort{NIR} oder \acrshort{NDVI}, und der Einbindung des räumlichen Kontexts erlaubt es, Objekte präziser zu erkennen und von Hintergrundstrukturen zu unterscheiden. Eine weitere Herausforderung ist die Wahl der Bounding Box-Strategie: \Acrfullpl{ABB} können bei schräg ausgerichteten oder überlappenden Objekten ungenau sein, während \Acrfullpl{OBB} eine präzisere Lokalisierung und Klassifikation kleiner Objekte ermöglichen. Diese Fähigkeit ist bei klein dargestellten Objekten in heterogenen Landschaften, wie sie in Luft- oder Satellitenbildern vorkommen, von zentraler Bedeutung für eine zuverlässige Detektion und Klassifikation.

Es ist dabei zu berücksichtigen, dass die ursprüngliche \acrfull{YOLO}-Architektur nicht speziell für die Analyse von Fernerkundungsdaten entwickelt wurde, sondern primär für die Objekterkennung in alltäglichen Szenen, wie sie in konventionellen Bilddatensätzen (z.\,B. \acrshort{COCO} oder ImageNet) vorkommen. Entsprechend basiert das Modell auf einem anderen räumlich-sensorischen Kontext, der sich von Fernerkundungsanwendungen deutlich unterscheidet. Dieses sogenannte \emph{Domain Shift}-Problem führt dazu, dass die Leistungsfähigkeit von YOLO bei der Erkennung kleiner, aus der Vogelperspektive dargestellter Objekte eingeschränkt ist. Diese Modelle sind daher auf die Erkennung großer, klar abgegrenzter Objekte in einer typischen Bodenperspektive optimiert. Im Kontext der Fernerkundung ergeben sich jedoch grundlegend andere Bedingungen: Objekte erscheinen aufgrund der Aufnahmeperspektive deutlich kleiner, überlappen teilweise und sind in heterogenen, komplexen Landschaften eingebettet. 

Die daraus resultierende methodische Herausforderung besteht darin, dass klassische Detektionsmodelle wie YOLO den räumlichen Kontext von Objekten bislang kaum berücksichtigen. Ein zentrales Merkmal dieser Arbeit ist deshalb die explizite Einbindung des \emph{räumlichen Kontexts} in die Objekterkennung. Echte Objekte treten stets in realen, räumlich strukturierten Umgebungen auf, deren Hintergrund – beispielsweise Asphalt, Sand oder Vegetation – maßgeblich die Detektionsleistung beeinflusst. Die Integration solcher kontextuellen Informationen (bspw. durch \acrshort{NDVI}) in die Analyse stellt einen wesentlichen methodischen Schwerpunkt dieser Arbeit dar, da sie entscheidend zur Verbesserung der Erkennungsgenauigkeit in realistischen Fernerkundungsszenarien beiträgt.

Die vorliegende Masterarbeit knüpft an diese Ergebnisse an und überträgt die Fragestellung auf den Bereich der Detektion und Klassifikation von kleinen Fahrzeugen in hochauflösenden multispektralen Luftbilddaten. Dabei wird eine neuere Version von \acrfull{YOLO} (v9) eingesetzt, die Verbesserungen in der Netzwerkarchitektur aufweist. Zusätzlich wird untersucht, inwieweit die Integration elektro-optischer und multispektraler Luftbildaufnahmen die Präzision bei der Detektion und Klassifikation sehr kleiner Objekte steigern kann. Auf diese Weise erweitert die Arbeit bestehende Forschung sowohl inhaltlich (Fokus auf Fahrzeuge) als auch methodisch (Verwendung moderner YOLO-Architektur und multispektraler Daten) und adressiert die Forschungslücke der Objekterkennung in komplexen Szenarien. Die Erweiterung der Analyse um den räumlichen Kontext – insbesondere durch die Berücksichtigung verschiedener Hintergrundtypen wie Asphalt oder Sand – ermöglicht eine realitätsnähere Bewertung der Modellleistung und adressiert damit eine zentrale Herausforderung der Objekterkennung in Fernerkundungsbildern.



\section{Structure of the thesis}


Diese Thesis bietet einen umfassenden Überblick über die Methoden, Experimente und Ergebnisse der Forschung. Sie beginnt mit einer Einleitung, die die Motivation und die Ziele der Arbeit darlegt, gefolgt von den Grundlagen, in denen der theoretische Hintergrund und relevante Metriken erläutert werden. Der Stand der Forschung fasst die aktuelle Literatur zusammen, bevor die Methodik die verwendeten Datensätze, Werkzeuge und die Vorgehensweise im Detail beschreibt. Anschließend werden die Ergebnisse der Experimente präsentiert und in der Diskussion interpretiert und in den Kontext der bestehenden Forschung eingeordnet. Die Arbeit schließt mit einer Schlussfolgerung und einem Ausblick auf mögliche zukünftige Forschungsrichtungen. Ergänzendes Material wie Bilder und Tabellen findet sich im Anhang.